\section{LA PROPAGAZIONE DEI CAMPI}

\subsection{L'equazione delle onde}

In questo capitolo, lungi dal presentare una trattazione esauriente del problema
generale della determinazione dei campi note le sorgenti, cercheremo di capire,
soprattutto con l'ausilio di qualche semplice esempio, l'importantissimo
fenomeno della propagazione dei campi elettromagnetici (\textit{onde
elettromagnetiche}), e della loro produzione.\\
Il punto di partenza sono (ovviamente) le equazioni di Maxwell (6.53);
supponiamo di essere in una regione di spazio priva di sorgenti: le (6.53) si
scrivono allora
\begin{equation}
  \left\{\begin{aligned}
    &\dive\vec E=0\\[4pt]
    &\rot\vec E+\frac{1}{c}\frac{\partial\vec B}{\partial t}=0\\[4pt]
    &\dive\vec B=0\\[4pt]
    &\rot\vec B=\frac{1}{c}\frac{\partial\vec E}{\partial t}
  \end{aligned}\right.
  \tag{12.1}
\end{equation}
(equazioni di Maxwell ``nel vuoto''). Le (12.1) sono tutto sommato abbastanza
semplici, in quanto si tratta di un sistema lineare ed omogeneo di equazioni
differenziali, anche se l'insieme delle soluzioni (cio\`e la \textit{soluzione
generale}) costituisce uno spazio vettoriale a infinite dimensioni: la vera
difficolt\`a nasce quando fra queste infinite soluzioni occorre individuare
quella determinata dalle sorgenti, che stanno fuori della regione considerata.
Limitiamoci allora a trovare alcune (semplici) soluzioni delle (12.1), poi
cercheremo di capire quali sorgenti possono determinarle.\\
Conviene riscrivere le (12.1) in modo diverso, prendendo la rotazione di ambo i
membri della 2$^\mathrm{a}$ e della 4$^\mathrm{a}$ ed utilizzando l'identit\`a
(7.20) che, essendo (nel vuoto) sia $\vec E$ che $\vec B$ a divergenza nulla, si
scrive
\begin{equation}
  \rot(\rot\vec E\,)=-\Delta\,\vec E\,;\quad\ \rot(\rot\vec B\,)=-\Delta\,\vec B
  \tag{12.2}
\end{equation}
e quindi
\begin{equation}
  \begin{aligned}
    &\Delta\,\vec E-\frac{1}{c^2}\frac{\partial^2\vec E}{\partial t^2}=0\\
    &\Delta\,\vec B-\frac{1}{c^2}\frac{\partial^2\vec B}{\partial t^2}=0
  \end{aligned}
  \tag{12.3}
\end{equation}
che, rispetto alle (12.1), hanno il vantaggio di essere disaccoppiate. Le (12.3)
non sono equivalenti alle (12.1) \big(non contengono l'informazione che $\vec E$
e $\vec B$ sono a divergenza nulla, e neppure il legame fra $\vec E$ e $\vec B$
stabilito dalle (12.1)\big), tuttavia sono importanti in quanto, come vedremo,
determinano la propagazione ondosa dei campi, e per questa ragione equazioni
della forma (12.3) sono note come \textbf{equazioni delle onde}.\\
Ogni componente del campo elettrico e di quello magnetico obbedisce quindi
all'equazione delle onde
\begin{equation}
  \Delta\,f-\frac{1}{c^2}\frac{\partial^2f}{\partial t^2}=0
  \tag{12.4}
\end{equation}
nota anche come \textbf{equazione di d'Alembert}; in analogia con l'equazione di
Laplace, la (12.4) spesso viene scritta introducendo l'operatore $\Box$ \ di
d'Alembert:
\begin{equation}
  \Box\,f=0\,;\qquad\quad\Box\equiv\Delta
  -\frac{1}{c^2}\frac{\partial^2}{\partial t^2}\,\cdot
  \tag{12.4'}
\end{equation}

\subsection{Le onde piane}

Per capire il fenomeno della propagazione dei campi, consideriamo l'equazione
(12.4) e riduciamoci al \underline{caso unidimensionale}, andando a cercare
quelle soluzioni che dipendono da una sola variabile spaziale ($x$), oltre che
dal tempo: \ $f=f(x,t)$; queste soluzioni sono dette \textbf{onde piane}, in
quanto per ogni $t$ sono costanti su ogni piano $x=$ costante. In questo caso la
(12.4) diventa
\begin{equation}
  \frac{\partial^2f(x,t)}{\partial x^2}
  -\frac{1}{c^2}\frac{\partial^2f(x,t)}{\partial t^2}=0.
  \tag{12.5}
\end{equation}
Per risolvere la (12.5) conviene effettuare il cambiamento di variabili
\begin{equation}
  \left\{\begin{aligned}
    \xi&=x-ct\\
    \eta&=x+ct\,.
  \end{aligned}\right.
  \tag{12.6}
\end{equation}
Poich\'e
\[
  \begin{aligned}
    &\frac{\partial f}{\partial x}
    =\frac{\partial f}{\partial\xi}\frac{\partial\xi}{\partial x}
    +\frac{\partial f}{\partial\eta}\frac{\partial\eta}{\partial x}
    =\frac{\partial f}{\partial\xi}+\frac{\partial f}{\partial\eta}\,,
    \qquad\frac{1}{c}\frac{\partial f}{\partial t}
    =\frac{\partial f}{\partial\xi}-\frac{\partial f}{\partial\eta}
    \qquad\Rightarrow\\
    &\frac{\partial^2f}{\partial x^2}=\frac{\partial^2f}{\partial\xi^2}
    +\frac{\partial^2f}{\partial\xi\partial\eta}
    +\frac{\partial^2f}{\partial\eta\partial\xi}
    +\frac{\partial^2f}{\partial\eta^2}\,,
    \qquad\frac{1}{c^2}\frac{\partial^2f}{\partial t^2}
    =\frac{\partial^2f}{\partial\xi^2}
    -2\frac{\partial^2f}{\partial\xi\partial\eta}
    +\frac{\partial^2f}{\partial\eta^2}
  \end{aligned}
\]
quindi in termini delle nuove variabili la (12.5) si scrive
\begin{equation}
  \frac{\partial^2f}{\partial\xi\,\partial\eta}=0
  \tag{12.7}
\end{equation}
la cui soluzione generale \`e
\begin{equation}
  f(\xi,\eta)=f_+(\xi)+f_-(\eta)\equiv f_+(x-ct)+f_-(x+ct)
  \tag{12.8}
\end{equation}
dove $f_-$ ed $f_+$ sono due funzioni \underline{arbitrarie}: infatti
\begin{equation}
  \frac{\partial}{\partial\xi}\Big(\frac{\partial f}{\partial\eta}\Big)=0
  \tag{12.9}
\end{equation}
comporta che \ $\partial f/\partial\eta$ \ \`e indipendente da $\xi$, cio\`e
\ $\partial f/\partial\eta=g(\eta)$, che integrata d\`a la (12.8), dove $f_-$
\`e una primitiva di $g$, e $f_+$ \`e la costante arbitraria (rispetto a
$\eta$), e quindi \`e un'arbitraria funzione di $\xi$.

Quindi, nel caso unidimensionale, la soluzione generale della (12.5) \`e la
sovrapposizione di un'arbitraria funzione $f_+$ che, avendo valore costante nei
punti in cui $x-ct$ \`e costante, trasla con velocit\`a $c$ nella direzione
positiva dell'asse $x$ mantenendo inalterata la propria forma (figura 12.1), e
di una funzione $f_-$ che trasla, sempre con velocit\`a $c$, nella direzione
negativa dell'asse $x$.
\fig{fig12_1}{12.1}
La (12.5), se $c$ \ \`e una opportuna costante, descrive per esempio il moto
delle (piccole) perturbazioni lungo una corda tesa (``equazione della corda
vibrante''): in questo caso \ $y\equiv f(x,t)$ \ \`e lo spostamento all'istante
t del generico punto $x$ della corda, in direzione ortogonale alla corda (``onde
trasversali''), cio\`e la forma della corda all'istante $t$, e
\ $c=\sqrt{\tau/\mu}$ \ ($\tau$ \`e la tensione della corda e $\mu$ la sua
densit\`a lineare di massa) \`e la velocit\`a di propagazione delle
perturbazioni lungo la corda.\\
Descrive altres\`i la propagazione unidimensionale del suono in un mezzo (aria,
acqua\dots): in questo caso $f(x,t)$ \`e (indifferentemente) la pressione nel
punto $x$ (all'istante t), oppure la densit\`a, o anche lo spostamento medio (in
direzione dell'asse $x$) delle particelle del fluido rispetto alla loro
posizione quando il mezzo \`e in quiete (``onde longitudinali'').

Torniamo alle (12.3) e cerchiamo soluzioni della forma
\begin{equation}
  \vec E=\vec E(x-ct),\quad\ \vec B=\vec B(x-ct)
  \tag{12.10}
\end{equation}
cio\`e \textit{onde piane progressive}, che siano anche soluzioni delle
(12.1).\\
Dalla 1$^\mathrm{a}$ e dalla 3$^\mathrm{a}$ delle (12.1) segue che -- a parte il
caso di campi costanti che non ci interessa -- sono nulle le componenti $x$ di
$\vec E$ e di $\vec B$, e quindi che le onde elettromagnetiche sono trasversali.
Dalla 2$^\mathrm{a}$ delle (12.1), tenuto conto che $\vec E$ e $\vec B$ non
dipendono da $y$ e da $z$, che $E_x=B_x=0$, e che per le onde progressive
\ $\partial f/\partial t=-c\,\partial f/\partial x$, si ha
\begin{equation}
  \frac{\partial E_z}{\partial x}=\frac{1}{c}\frac{\partial B_y}{\partial t}
  =-\frac{\partial B_y}{\partial x}\,;
  \qquad\qquad\frac{\partial E_y}{\partial x}
  =-\frac{1}{c}\frac{\partial B_z}{\partial t}=\frac{\partial B_z}{\partial x}
  \tag{12.11}
\end{equation}
(la 4$^\mathrm{a}$ delle (12.1) porta allo stesso risultato) da cui si vede che
le componenti $(E_z,B_y)$ sono indipendenti dalle componenti $(E_y,B_z)$.
Possiamo allora porre
\begin{equation}
  E_y=B_z=0
  \tag{12.12}
\end{equation}
per cui dalle (12.11) (sempre escludendo campi costanti) si ha
\begin{equation}
  \vec E=\big(0,\,0,\,f(x-ct)\big);\quad\ \vec B=\big(0,\,-f(x-ct),\,0\big).
  \tag{12.13}
\end{equation}
Una soluzione del tipo (12.13) \`e detta onda piana \textbf{polarizzata
linearmente} in direzione dell'asse $z$, in quanto il campo elettrico \`e in
ogni punto parallelo all'asse $z$.\\
Analogamente, se invece delle (12.12) poniamo \ $E_z=B_y=0$, si ha
\begin{equation}
  \vec E=\big(0,\,g(x-ct),\,0\big);\quad\ \vec B=\big(0,\,0,\,g(x-ct)\big).
  \tag{12.14}
\end{equation}
La soluzione (12.14) \`e un'onda piana polarizzata in direzione dell'asse $y$.\\
Poich\'e le (12.1) sono lineari e omogenee, anche qualsiasi combinazione lineare
delle (12.13) e delle (12.14) \`e soluzione delle (12.1):
\begin{equation}
  \vec E=\big(0,\,g(x-ct),\,f(x-ct)\big);
  \quad\ \vec B=\big(0,\,-f(x-ct),\,g(x-ct)\big).
  \tag{12.15}
\end{equation}
Si noti che $\vec E$ e $\vec B$ hanno lo stesso modulo e sono ortogonali:
\begin{equation}
  \vec E^2-\vec B^2=0\,;\quad\ \vec E\cdot\vec B=0
  \tag{12.16}
\end{equation}
cio\`e entrambi gli invarianti relativistici del campo elettromagnetico sono
nulli.\\
Per le onde che viaggiano in verso opposto all'asse $x$ la trattazione \`e
simile, e si ottiene
\begin{equation}
  \vec E=\big(0,\,g(x+ct),\,f(x+ct)\big);
  \quad\ \vec B=\big(0,\,f(x+ct),\,-g(x+ct)\big).
  \tag{12.17}
\end{equation}
Nel caso delle onde piane (progressive oppure regressive) il vettore di
Poynting, (quindi sia la corrente di energia che la densit\`a di quantit\`a di
moto) ha la direzione coincidente con la direzione di propagazione dell'onda ed
inoltre energia e quantit\`a di moto sono proporzionali; infatti dalle (12.15) e
dalle (11.25)
\begin{equation}
  \begin{aligned}
    &u(x,t)=\frac{1}{4\pi}\big(f^2(x-ct)+g^2(x-ct)\big);
    \qquad\vec S=(S,\,0,\,0),\quad S=cu(x,t)\quad\Rightarrow\\
    &U=cP\,.
  \end{aligned}
  \tag{12.18}
\end{equation}

\subsection{Dalle sorgenti ai campi: le soluzioni ritardate}

Le (12.15) e (12.17), anche se sono solo delle soluzioni particolari delle
equazioni di Maxwell, mettono in evidenza che le onde elettromagnetiche, e con
esse la densit\`a di energia e di quantit\`a di moto, si propagano con
velocit\`a $c$; ora vogliamo andare oltre, e mostrare quali sorgenti possono
produrre i campi descritti dalle (12.15) e (12.17); verificheremo anche che le
sorgenti in un dato punto $Q$ ad un dato istante contribuiscono ai campi in un
punto $P$ con un ritardo uguale al tempo $\overline{QP}\,/c$ che un segnale che
viaggia con velocit\`a $c$ impiega a propagarsi dal ``punto potenziante'' $Q$ al
``punto potenziato'' $P$ (il ritardo che viene trascurato quando si fa
l'approssimazione delle correnti lentamente variabili).\\
Siccome le onde piane che si propagano in direzione dell'asse $x$ sono
invarianti per traslazioni parallele al piano $x=0$, prendiamo sorgenti con le
stesse simmetrie: una lamina di corrente nel piano $x=0$
\begin{equation}
  \vec{\mathcal J}(t)=\vec{\mathcal J}_0\,f(t).
  \tag{12.19}
\end{equation}
Conviene calcolare il potenziale vettore $\vec A(x,y,z,t)$ la cui esistenza,
anche nel caso dei campi variabili, \`e garantita dal fatto che $\vec B$ \`e
solenoidale. Per dimostrare l'effetto ritardato delle sorgenti, calcoliamo il
potenziale vettore con formula analoga alla (7.22), nella quale per\`o la
sorgente nel punto potenziante $(x',y',z')$ contribuisce al campo all'istante
$t$ nel punto potenziato $(x,y,z)$ con il valore che ha all'istante $t-r/c$,
dove $r$ \`e la distanza fra punto potenziante e punto potenziato (e quindi i
vari punti potenzianti contribuiscono ad istanti diversi):
\begin{equation}
  \begin{aligned}
    &\vec A(x,y,z,t)=\frac{1}{c}
    \int_{x=0}\frac{\vec{\mathcal J}(t-r/c)}{r}\,dS\\
    &\Big(r\equiv\sqrt{(x-x')^2+(y-y')^2+(z-z')^2}\ \Big)
  \end{aligned}
  \tag{12.20}
\end{equation}
poi dimostreremo che i campi cos\`i determinati sono la soluzione del
problema.\\
Per calcolare il potenziale vettore in un (generico) punto $P$ effettuiamo
l'integrazione usando coordinate polari $\varrho,\phi$ nel piano delle sorgenti,
con l'origine coincidente con la proiezione di $P$ su detto piano; $P$ ha quindi
coordinate $(x,0,0)$ e $r=\sqrt{x^2+\rho^2}$:
\begin{equation}
  \vec A=\frac{\vec{\mathcal J}_0}{c}\int_0^\infty
  \frac{f\big(t-\sqrt{x^2+\varrho^2}/c\big)}{\sqrt{x^2+\varrho^2}}
  \,2\pi\varrho\,d\varrho\,.
  \tag{12.21}
\end{equation}
Posto
\begin{equation}
  \xi=\sqrt{x^2+\varrho^2}\quad\ \Rightarrow
  \quad\ \frac{\varrho}{\sqrt{x^2+\varrho^2}}=\frac{d\xi}{d\varrho}
  \tag{12.22}
\end{equation}
e quindi
\begin{equation}
  \vec A(x,t)
  =\frac{2\pi\vec{\mathcal J}_0}{c}\int_{|x|}^\infty f(t-\xi/c)\,d\xi\,.
  \tag{12.23}
\end{equation}
Se prendiamo l'asse $y$ nella direzione di $\vec{\mathcal J}_0$, si ha
\begin{equation}
  B_x=0,\ B_y=0,\ B_z=\frac{\partial A_y}{\partial x}=\left\{\begin{array}{ll}
    -\dfrac{2\pi\mathcal J_0}{c}f(t-x/c)&\quad x>0\\[10pt]
    \phantom{-}\dfrac{2\pi\mathcal J_0}{c}f(t+x/c)&\quad x<0
  \end{array}\right.
  \tag{12.24}
\end{equation}
e quindi dalla (12.14) (per l'onda progressiva) e dalla (12.17) (per l'onda
regressiva)
\begin{equation}
  E_x=0,\ E_z=0,\ E_y=\left\{\begin{array}{ll}
    -\dfrac{2\pi\mathcal J_0}{c}f(t-x/c)&\quad x>0\\[10pt]
    -\dfrac{2\pi\mathcal J_0}{c}f(t+x/c)&\quad x<0\,.
  \end{array}\right.
  \tag{12.25}
\end{equation}
Si noti che la lamina di corrente \`e una superficie di discontinuit\`a per
$\vec B$, ma non per $\vec E$.\\
Le (12.24) e (12.25) sono soluzioni per $x>0$ e per $x<0$; resta da dimostrare
che effettivamente costituiscono la soluzione corrispondente alle sorgenti date:
a questo scopo basta notare che il campo magnetico ha la corretta
discontinuit\`a \ $\Delta B_z=4\pi\mathcal J_0f(t)/c$ \ attraverso la lamina di
corrente (la corrente di spostamento non contribuisce alla discontinuit\`a).\\
La soluzione trovata consiste quindi di un'onda piana polarizzata nella
direzione della corrente, che si allontana (sia per $x>0$ che per $x<0$) dalle
sorgenti.\\
Vale la pena confrontare le (12.24), (12.25) con la soluzione calcolata
nell'approssimazione delle correnti lentamente variabili, che d\`a luogo a campi
indipendenti dalla distanza dal piano delle sorgenti
\ $\big(f(t\pm x/c)\to f(t)\big)$.

\subsection{Le onde piane \textit{monocromatiche}}

Le onde monocromatiche, non necessariamente piane, sono particolari onde che
dipendono dal tempo in modo periodico esclusivamente tramite le funzioni
$\sen\omega t$ e $\cos\omega t$:
\begin{equation}
  F(x,y,z,t)=a(x,y,z)\sen\omega t+b(x,y,z)\cos\omega t
  \tag{12.26}
\end{equation}
e quindi, siccome \ $\partial^2F/\partial t^2=-\omega^2F$, le funzioni spaziali
$a(x,y,z)$ e $b(x,y,z)$ \big(che genericamente indichiamo con $f(x,y,z)$\big)
soddisfano (nel vuoto) l'equazione
\begin{equation}
  \Delta\,f+\frac{\omega^2}{c^2}f=0.
  \tag{12.27}
\end{equation}
Siccome le onde piane sono funzioni di $(x\pm ct)$, le \textbf{onde piane
monocromatiche} sono combinazioni lineari di seni e coseni dell'argomento
$(\omega/c)(x\pm ct)$, e quindi sono tutte della forma
\begin{equation}
  F(x,t)=f_0\cos\big(\frac{\omega}{c}(x\pm ct)+\phi\big)
  \equiv f_0\cos\big(kx\pm\omega t+\phi\big);\qquad k=\frac{\omega}{c}
  \tag{12.28}
\end{equation}
perci\`o esse non solo sono, per ogni fissato $x$, periodiche nel tempo con
\textbf{periodo} \ $T=2\pi/\omega$, e quindi con \textbf{frequenza}
\ $\nu=\omega/2\pi$, ma, per ogni $t$ fissato, sono funzioni periodiche in $x$
con \textbf{lunghezza d'onda}, cio\`e periodo spaziale,
\ $\lambda=2\pi/k=2\pi c/\omega=c/\nu$:
\begin{equation}
  T=\frac{2\pi}{\omega}\,,\quad\ \nu=\frac{1}{T}=\frac{\omega}{2\pi}\,,
  \quad\ \lambda=\frac{c}{\nu}\,,\quad\ \nu=\frac{c}{\lambda}\,\cdot
  \tag{12.29}
\end{equation}
La fase $\phi$ nella (12.28) non ha un significato fisico particolare, in quanto
dipende dalla scelta dell'origine delle coordinate e/o dei tempi.\\
Pi\`u in generale, se $\hat n$ \`e la direzione di propagazione dell'onda,
possiamo scrivere
\begin{equation}
  F(x,y,z,t)=f_0\cos\big(\vec k\cdot\vec r-\omega t+\phi\big),
  \qquad\vec k\equiv k\hat n
  \tag{12.30}
\end{equation}
ed il vettore $\vec k$ \`e detto \textbf{vettore d'onda}.\\
Le onde elettromagnetiche piane monocromatiche (progressive in direzione
dell'asse $x$) sono della forma
\begin{equation}
  \vec E(x,t)
  =\big(0,\,E_{0y}\cos(kx-\omega t+\phi),\,E_{0z}\cos(kx-\omega t+\psi)\big).
  \tag{12.31}
\end{equation}
Grazie alle (12.15) \`e sufficiente l'espressione del campo elettrico: $\vec B$
\`e univocamente determinato da $\vec E$ e, come in ogni onda piana, $\vec E$ ed
$\vec B$ hanno lo stesso modulo ed insieme alla direzione di propagazione
formano (nell'ordine) una terna ortogonale destrorsa. Per quanto detto sopra,
$\phi$ e $\psi$ nella (12.31) non hanno separatamente significato fisico, ma lo
ha solo la loro differenza.\\
L'\textbf{intensit\`a} $I$ dell'onda \`e definita come l'energia (media su un
periodo) che attraversa nell'unit\`a di tempo l'unit\`a di superficie ortogonale
alla direzione di propagazione; essa \`e quindi data dal modulo del vettore di
Poynting mediato su un periodo; grazie al fatto che $\vec E$ e $\vec B$ sono
ortogonali e hanno lo stesso modulo si ha
\begin{equation}
  I\equiv\frac{1}{T}\int_t^{t+T}|\,\vec S\,|\,dt'
  =\frac{1}{T}\int_t^{t+T}\frac{c}{4\pi}\vec E^2\,dt'=\frac{c}{8\pi}\vec E_0^2;
  \quad\ \vec E_0\equiv(0,\,E_{0y},\,E_{0z})
  \tag{12.32}
\end{equation}
(la media di $\cos^2(\omega t+\phi)$ \`e 1/2).\\
La traiettoria descritta dal vettore $\vec E$ in ogni fissato punto
$\overline P$ dello spazio \`e data dalle equazioni parametriche
\begin{equation}
  \left\{\begin{aligned}
    &E_x(t)=0\\
    &E_y(t)=E_{0y}\cos(k\overline x-\omega t+\phi)\\
    &E_z(t)=E_{0z}\cos(k\overline x-\omega t+\psi)
  \end{aligned}\right.
  \tag{12.33}
\end{equation}
che sono (in generale) le equazioni di un'ellisse, che viene interamente
descritta in un periodo: se $\phi=\psi$ l'ellisse degenera in un segmento,
cio\`e il campo elettrico (in ogni punto) oscilla nella direzione (parallela al
piano $x=0$) che forma con l'asse $y$ l'angolo $\theta$ dato da
\ $\tan\theta=E_{0z}/E_{0y}$, quindi l'onda \`e polarizzata linearmente (onda
$\pi$), ed in questo caso la (12.31) assume la forma
\begin{equation}
  \vec E(x,t)=\vec E_0\cos(kx-\omega t+\phi),\qquad(\vec E_0\cdot\vec
  k=0)\qquad\qquad\text{(onda $\pi$)}.
  \tag{12.34}
\end{equation}
Se \ $E_{0y}=E_{0z}$ \ e \ $\phi-\psi=\pm\pi/2$, l'ellisse si riduce ad un
cerchio, percorso in verso orario (per un osservatore che vede l'onda propagarsi
verso di lui) se \ $\phi-\psi=\pi/2$, e in verso antiorario se
\ $\phi-\psi=-\pi/2$: nel primo caso si parla di polarizzazione circolare destra
(onda $\sigma_+$), e nel secondo di polarizzazione circolare sinistra (onda
$\sigma_-$):
\begin{equation}
  \vec E(x,t)=\frac{E_0}{\sqrt2}\big(0,\,\mp\sen(kx-\omega
  t+\phi),\,\cos(kx-\omega t+\phi)\big)\qquad\text{(onde $\sigma_\pm$)}.
  \tag{12.35}
\end{equation}
In tutti gli altri casi si dice che l'onda \`e polarizzata ellitticamente.
Le onde piane monocromatiche sono quindi caratterizzate dalla frequenza (o
equivalentemente dalla lunghezza d'onda), dall'intensit\`a, dalla direzione di
propagazione e dalla polarizzazione.\\
Le onde radio hanno lunghezze d'onda dell'ordine di 10\,m (onde corte), 100\,m
(onde medie), 1000\,m (onde lunghe), mentre per le trasmissioni in modulazione
di frequenza e per la televisione vengono utilizzate onde con frequenze attorno
ai 100\,Mh \ ($\lambda\approx3$\,m); il radar utilizza onde centimetriche
\ ($\nu\approx10^{10}$\,Hz); attorno ai $10^{14}$\,Hz abbiamo la radiazione
infrarossa; la luce visibile ha lunghezza d'onda fra i 7000\,\AA{} e i
3000\,\AA{}. Pi\`u oltre (lunghezza d'onda minore) la luce ultravioletta, i
raggi $X$ \ ($\nu\approx10^{18}$\,Hz) ed infine i raggi $\gamma$.

Le onde piane sono ovviamente non fisiche, in quanto richiedono sorgenti
distribuite in regioni di estensione infinita; a maggior ragione sono non
fisiche le onde piane monocromatiche, che sono non nulle per ogni $t$, da
$t=-\infty$ a $t=+\infty$ e quindi la stessa cosa deve valere per le sorgenti.\\
Tuttavia, da un lato la loro non--fisicit\`a non \`e diversa da quella dei
solenoidi infiniti, dei fili rettilinei indefiniti \dots, e persino da quella
dei campi statici che richiedono che la posizione delle sorgenti sia immutabile
da un tempo infinito: esse quindi possono rappresentare una buona
approssimazione in regioni limitate dello spazio lontane dalle sorgenti, e per
intervalli di tempo finiti; d'altro lato, grazie alla linearit\`a delle
equazioni di Maxwell nel vuoto, la somma (finita, infinita anche con la potenza
del continuo) di onde piane di diversa frequenza, intensit\`a, direzione di
propagazione e polarizzazione, \`e ancora soluzione delle equazioni di Maxwell,
ed \`e possibile dimostrare che in questo modo si pu\`o ottenere qualsiasi
soluzione.
\fig{fig12_2}{12.2}

\subsection{Un caso realistico: l'onda emessa da un'antenna}

Una antenna \`e sostanzialmente un conduttore, che assumiamo rettilineo e di
lunghezza $2l$, nel quale un generatore mantiene una corrente variabile (figura
12.2). Presa l'origine nel centro dell'antenna, e l'asse $z$ lungo la stessa,
supponiamo che la corrente sia della forma
\begin{equation}
  I(z,t)=I_0(z)f(t).
  \tag{12.36}
\end{equation}
$I_0(z)$ si annulla agli estremi dell'antenna, e per esempio pu\`o essere della
forma \ $I_0(z)=I_0\sen k(l-|z|)$.\\
Vogliamo determinare i campi a grande distanza: a questo scopo calcoleremo il
potenziale vettore, prendendo per esso l'espressione ``ritardata'' simile alla
(12.20)
\begin{equation}
  A_z(\vec r,t)=\frac{1}{c}\int_{-l}^l\frac{I(z',t-|\vec r-\vec r\,'|/c)}
  {|\vec r-\vec r\,'|}\,dz'
  \tag{12.37}
\end{equation}
che non dimostriamo, ma che per lo meno abbiamo visto funzionare bene in un caso
particolare (siccome la corrente \`e parallela all'asse $z$, le altre componenti
di $\vec A$ sono nulle).\\
Oltre all'approssimazione delle grandi distanze, che consiste nel sostituire
$|\vec r-\vec r\,'|$ con
$r$ al denominatore della (12.37), faremo un'altra approssimazione, nota come
\textbf{approssimazione di dipolo}: poniamo
\begin{equation}
  f(t-|\vec r-\vec r\,'|/c)\approx f(t-r/c).
  \tag{12.38}
\end{equation}
Questa approssimazione richiede non solo che $r\gg l$, ma anche che la corrente
vari poco nell'intervallo di tempo $l/c$ \ che la luce impiega a percorrere la
distanza uguale alla lunghezza dell'antenna. Infatti, grazie allo sviluppo
(1.75),
\begin{equation}
  \begin{aligned}
    f(t-|\vec r-\vec r\,'|/c)&\approx f(t-r/c+\hat r\cdot\vec r\,'/c)\\
    &\approx f(t-r/c)+\frac{1}{c}\dot f(t-r/c)\,\hat r\cdot\vec r\,'
  \end{aligned}
  \tag{12.39}
\end{equation}
e quindi deve essere
\begin{equation}
  \dot f\cdot\frac{l}{c}\ll f.
  \tag{12.40}
\end{equation}
Questa \`e simile all'approssimazione delle correnti lentamente variabili,
per\`o \`e una richiesta molto pi\`u debole, in quanto ristretta solo alla
regione dell'antenna, mentre invece i campi verranno calcolati a grande
distanza, e quindi in una regione molto pi\`u ampia. Se -- per esempio -- la
corrente nell'antenna \`e sinusoidale, cio\`e \ $f(t)=\cos\omega t$, la (12.40)
equivale a \ $\omega l/c\ll1$, cio\`e \ $\lambda\gg l$.\\
In definitiva, grazie alle approssimazioni fatte e alla (12.36), la (12.37)
diventa
\begin{equation}
  A_z(\vec r,t)=\frac{1}{rc}f(t-r/c)\int_{-l}^lI_0(z')\,dz'.
  \tag{12.41}
\end{equation}
L'integrale nella (12.41) \`e una costante che ha le dimensioni della derivata
rispetto al tempo di una carica per una lunghezza, cio\`e della derivata
rispetto al tempo di un momento di dipolo, per cui indicheremo il prodotto di
questa costante con $f(t)$ con $\dot p(t)$:
\begin{equation}
  A_z(\vec r,t)=\frac{\dot p(t-r/c)}{rc}\,\cdot
  \tag{12.42}
\end{equation}
Si noti che a grandi distanze il potenziale vettore \`e di ordine $1/r$, e
\underline{non} $1/r^2$ \ come nel caso statico \big((7.50)\big): la ragione sta
nel fatto che ora la corrente non \`e a divergenza nulla, per cui il suo
integrale esteso al circuito non \`e nullo.\\
Per calcolare $\vec B$ dobbiamo prendere la rotazione di $\vec A$: nel calcolo
dobbiamo trascurare i termini $O(1/r^2)$ che provengono dalle derivate di $1/r$
nella (12.42), quindi il contributo al campo magnetico a grandi distanze
proviene solo dalle derivate del numeratore nella (12.42); poich\'e
\ $\partial r/\partial x_i=x_i/r$, si ha
\begin{equation}
  \left\{\begin{aligned}
    B_x&=\phantom{-}\frac{\partial A_z}{\partial y}
    =-\frac{\ddot p(t-r/c)}{rc^2}\,\frac{y}{r}\\
    B_y&=-\frac{\partial A_z}{\partial x}
    =\phantom{-}\frac{\ddot p(t-r/c)}{rc^2}\,\frac{x}{r}\\
    B_z&=\phantom{-}0
  \end{aligned}\right.
  \tag{12.43}
\end{equation}
che pu\`o essere scritta nella forma pi\`u compatta
\[
  \vec B(\vec r,t)=-\frac{\ddot p(t-r/c)}{rc^2}\,\hat r\wedge\hat\kappa
\]
o anche, posto \ $\vec p\equiv p\,\hat\kappa$,
\begin{equation}
  \vec B(\vec r,t)=\frac{1}{rc^2}\,\ddot{\vec p}(t-r/c)\wedge\hat r\,.
  \tag{12.44}
\end{equation}
Il campo elettrico si calcola dall'equazione
\ $\partial\vec E/\partial t=c\rot\vec B$, al solito trascurando i termini
$O(1/r^2)$, ed escludendo la presenza di un campo statico:
\begin{equation}
  \vec E(\vec r,t)=\frac{1}{rc^2}\,\big(\ddot{\vec p}(t-r/c)\wedge\hat r\big)
  \wedge\hat r\,.
  \tag{12.45}
\end{equation}
Ci sono molte cose da dire a proposito del risultato ottenuto, alcune delle
quali hanno carattere generale:
\begin{enumerate}
  \item i campi dovuti a sorgenti variabili, contenute in una regione limitata
    dello spazio, a grande distanza dalle sorgenti hanno caratteristiche diverse
    dai campi generati da sorgenti statiche (o stazionarie): questi ultimi
    decrescono almeno come $1/r^2$ nel caso del campo elettrico, e come $1/r^3$
    nel caso di quello magnetico; lo sviluppo dei campi variabili in potenze
    inverse della distanza dalle sorgenti, invece, comincia (o meglio: pu\`o
    cominciare) con termini di ordine $1/r$, sia nel caso del campo elettrico
    che di quello magnetico. I termini di ordine $1/r$ costituiscono il
    \textbf{campo di radiazione};
  \item tanto $\vec E$ quanto $\vec B$ dati dalle (12.44) e (12.45) sono (dove
    non nulli) ortogonali fra di loro, hanno lo stesso modulo e sono tangenti
    alle sfere centrate sull'antenna, cio\`e sono ortogonali al vettore $\vec r$
    dall'antenna al punto di osservazione (o ``punto potenziato''); $\vec E$ \`e
    diretto lungo i meridiani mentre il campo magnetico \`e diretto lungo i
    paralleli: come per le onde piane, $\vec E,\vec B$ e la direzione di
    osservazione $\hat r$ formano una terna destrorsa (figura 12.3); lungo ogni
    direzione di osservazione l'onda \`e polarizzata linearmente;
    \fig{fig12_3}{12.3}
  \item il vettore di Poynting, cio\`e la corrente di energia, \`e parallelo ad
    $\vec r$ e quindi le linee di forza di $\vec S$ sono rette: ci\`o si esprime
    dicendo che la propagazione delle onde elettromagnetiche (nel vuoto) \`e
    \textbf{rettilinea}.
\end{enumerate}
Supponiamo ora che la corrente nell'antenna sia sinusoidale; in questo caso si
ha \ $p(t)=p_0\sen\omega t$, e quindi
\begin{equation}
  \left\{\begin{aligned}
    &\vec E(\vec r,t)=\frac{\omega^2}{rc^2}
    \sen(kr-\omega t)\,(\vec p_0\wedge\hat r)\wedge\hat r\\
    &\vec B(\vec r,t)=\frac{\omega^2}{rc^2}
    \sen(kr-\omega t)\,\vec p_0\wedge\hat r
  \end{aligned}\right.
  \tag{12.46}
\end{equation}
quindi l'onda emessa dall'antenna \`e monocromatica (nell'approssimazione di
dipolo).\\
Siccome ($\theta$ \`e l'angolo fra $\vec r$ e $\vec p_0$)
\begin{equation}
  |\vec p_0\wedge\hat r|=|(\vec p_0\wedge\hat r)\wedge\hat r|=p_0\sen\theta
  \tag{12.47}
\end{equation}
si ha
\begin{equation}
  \begin{aligned}
    &\vec S\equiv\frac{c}{4\pi}\vec E\wedge\vec B
    =\frac{c}{4\pi}E\,B\,\hat r=\frac{p_0^2\omega^4\sen^2\theta}{4\pi r^2c^3}\,
    \sen^2(kr-\omega t)\,\hat r,\\
    &I(\vec r\,)=\overline{S(\vec r,t)}
    =\frac{p_0^2\omega^4\sen^2\theta}{8\pi r^2c^3}
  \end{aligned}
  \tag{12.48}
\end{equation}
quindi l'emissione di energia da parte dell'antenna non \`e isotropa; \`e
massima in direzione ortogonale all'antenna ed \`e nulla nella direzione
parallela: un ricevitore allineato con l'antenna non riceve alcun segnale;
inoltre, in una data posizione, la ricezione del segnale \`e migliore se
l'antenna del ricevitore \`e parallela alla direzione di polarizzazione
dell'onda.\\
Nella figura 12.4 \`e riportata una rappresentazione grafica dell'intensit\`a in
funzione dell'angolo $\theta$ fra la direzione di osservazione e quella
dell'antenna: $I(r,\theta)$ \`e proporzionale alla lunghezza del segmento fra
l'origine e il punto d'incontro con la curva, sulla semiretta che forma un
angolo $\theta$ con l'antenna.
\fig{fig12_4}{12.4}
L'energia media (su un periodo) che nell'unit\`a di tempo attraversa una sfera
di raggio $r$ per la (12.48) \`e
\begin{equation}
  W=\frac{p_0^2\omega^4}{8\pi r^2c^3}2\pi r^2
  \int_{-1}^1\sen^2\theta\,d\cos\theta=\frac{p_0^2\omega^4}{3c^3}
  \tag{12.49}
\end{equation}
ed \`e indipendente dal raggio (per $r\gg l$); ci\`o \`e dovuto al fatto che,
essendo il campo di radiazione di ordine $1/r$, la corrente di energia \`e di
ordine $1/r^2$, mentre la superficie della sfera \`e proporzionale a $r^2$: \`e
quindi il campo di radiazione che \`e responsabile del trasporto di energia in
regioni comunque lontane dalle sorgenti. Quindi la (12.49) \`e anche la potenza
media che il generatore deve erogare all'antenna.

Le (12.44) e (12.45) valgono in tutti i casi in cui si ha un sistema di cariche
il cui momento di dipolo $\vec p$ varia nel tempo, con derivata seconda diversa
da zero: per esempio una carica oscillante, o pi\`u in generale una carica
accelerata. L'irraggiamento quindi richiede che le cariche abbiano
un'accelerazione: in effetti sappiamo che una carica in moto uniforme genera un
campo che decresce come $1/r^2$, e quindi non \`e un campo di radiazione.

\paragraph{Esempio}
\begin{quote}
  Calcoliamo la potenza media irraggiata da un'antenna alimentata da un
  generatore che eroga una corrente sinusoidale di pulsazione
  \ $\omega=10^8\,\mathrm{s}^{-1}$ \ tale che
  \begin{equation}
    \int_{-l}^lI_0(z)\,dz=10\,\mathrm{A\,m}
    \tag{12.50}
  \end{equation}
  (per esempio un'antenna di 1\,m alimentata con una corrente media di 10\,A).\\
  Siccome
  \begin{equation}
    p(t)=p_0\sen\omega t\ \Rightarrow\ \dot p(t)=p_0\omega\cos\omega t
    \ \Rightarrow\ p_0\omega=\int_{-l}^lI_0(z)\,dz
    \tag{12.51}
  \end{equation}
  e dalla (12.49) si ottiene
  \begin{equation}
    W=\frac{(p_0\omega)^2\omega^2}{3c^3}
    =\frac{(3\times10^{10})^2\cdot10^4\cdot10^{16}}{3\cdot(3\times10^{10})^3}
    \simeq10^9\,\mathrm{erg}=100\,\mathrm{watt}.
    \tag{12.52}
  \end{equation}
  La lunghezza d'onda della radiazione emessa \`e
  \begin{equation}
    \lambda=\frac{2\pi c}{\omega}=20\,\mathrm{m}
    \tag{12.53}
  \end{equation}
  quindi se la lunghezza dell'antenna \`e dell'ordine del metro,
  l'approssimazione di dipolo pu\`o essere considerata accettabile.
\end{quote}

Infine, vediamo esplicitamente in quali condizioni la soluzione data dalle
(12.46) pu\`o essere approssimata con un'onda piana: $r$ al denominatore delle
(12.46) pu\`o essere considerato costante ($r=r_0$) se all'interno della regione
di osservazione esso varia poco rispetto alla distanza dalle sorgenti:
\ $|r-r_0|\ll r_0$; prendiamo l'asse $x$ diretto dalle sorgenti alla regione di
osservazione, e sia $d$ la dimensione trasversa all'asse $x$ di detta regione.
All'interno della regione di osservazione ($x\approx r_0$), e se $d\ll r_0$, si
ha
\begin{equation}
  r\equiv\sqrt{x^2+y^2+z^2}\approx x\Big(1+\frac{y^2+z^2}{2x^2}\Big)
  \tag{12.54}
\end{equation}
per cui \ $kr\approx kx\big(1+(y^2+z^2)/2x^2\big)$ \ e quindi, essendo
\ $d^2\ll r_0^2$, possiamo porre nella (12.46)
\begin{equation}
  \frac{\sen(kr-\omega t)}{r}\approx\frac{\sen(kx-\omega t)}{r_0}\,\cdot
  \tag{12.55}
\end{equation}
In definitiva, come anticipato, se le dimensioni sia longitudinali ($\Delta x$)
che trasverse ($d$) della regione di osservazione sono molto minori della
distanza dalle sorgenti, la curvatura dei fronti d'onda e l'attenuazione con la
distanza possono essere trascurate e l'onda descritta dalla (12.46) pu\`o essere
approssimata con un'onda piana.
